\documentclass{article}
\usepackage[utf8]{inputenc}
\usepackage[T1]{fontenc}
\usepackage{geometry}
\usepackage{graphicx}
\usepackage{float}
\usepackage{hyperref}
\usepackage[english,indonesian]{babel}

\title{Optimasi Pemrosesan Template LaTeX Menggunakan Pendekatan Berbasis AI}
\author{Academic Support Agent}
\date{\today}

\begin{document}
\maketitle

\begin{abstract}

\end{abstract}

\chapter{PENDAHULUAN}


\chapter{TINJAUAN PUSTAKA}


% Original content (commented out):
% % Options for packages loaded elsewhere
\PassOptionsToPackage{unicode}{hyperref}
\PassOptionsToPackage{hyphens}{url}
%
\documentclass[11pt
]{article}
\usepackage{amsmath,amssymb}
\usepackage{lmodern}
\usepackage{iftex}
\ifPDFTeX
  \usepackage[T1]{fontenc}
  \usepackage[utf8]{inputenc}
  \usepackage{textcomp} % provide euro and other symbols
\else % if luatex or xetex
  \usepackage{unicode-math}
  \defaultfontfeatures{Scale=MatchLowercase}
  \defaultfontfeatures[\rmfamily]{Ligatures=TeX,Scale=1}
\fi
% Use upquote if available, for straight quotes in verbatim environments
\IfFileExists{upquote.sty}{\usepackage{upquote}}{}
\IfFileExists{microtype.sty}{% use microtype if available
  \usepackage[]{microtype}
  \UseMicrotypeSet[protrusion]{basicmath} % disable protrusion for tt fonts
}{}
\makeatletter
\@ifundefined{KOMAClassName}{% if non-KOMA class
  \IfFileExists{parskip.sty}{%
    \usepackage{parskip}
  }{% else
    \setlength{\parindent}{0pt}
    \setlength{\parskip}{6pt plus 2pt minus 1pt}}
}{% if KOMA class
  \KOMAoptions{parskip=half}}
\makeatother
\usepackage{xcolor}
\IfFileExists{xurl.sty}{\usepackage{xurl}}{} % add URL line breaks if available
\IfFileExists{bookmark.sty}{\usepackage{bookmark}}{\usepackage{hyperref}}
\hypersetup{
  pdftitle={An Example Abstract for the International Conference on Electron, Ion, and Photon Beam Technology and Nanofabrication},
  hidelinks,
  pdfcreator={LaTeX via pandoc}}
\urlstyle{same} % disable monospaced font for URLs
\usepackage{graphicx}
\makeatletter
\def\maxwidth{\ifdim\Gin@nat@width>\linewidth\linewidth\else\Gin@nat@width\fi}
\def\maxheight{\ifdim\Gin@nat@height>\textheight\textheight\else\Gin@nat@height\fi}
\makeatother
% Scale images if necessary, so that they will not overflow the page
% margins by default, and it is still possible to overwrite the defaults
% using explicit options in \includegraphics[width, height, ...]{}
\setkeys{Gin}{width=\maxwidth,height=\maxheight,keepaspectratio}
% Set default figure placement to htbp
\makeatletter
\def\fps@figure{htbp}
\makeatother
\usepackage[normalem]{ulem}
% Avoid problems with \sout in headers with hyperref
\pdfstringdefDisableCommands{\renewcommand{\sout}{}}
\setlength{\emergencystretch}{3em} % prevent overfull lines
\providecommand{\tightlist}{%
  \setlength{\itemsep}{0pt}\setlength{\parskip}{0pt}}
\setcounter{secnumdepth}{-\maxdimen} % remove section numbering
\ifLuaTeX
  \usepackage{selnolig}  % disable illegal ligatures
\fi

\usepackage[margin=3cm]{geometry}

\title{An Example Abstract for the International Conference on Electron,
Ion, and Photon Beam Technology and Nanofabrication}
\usepackage{etoolbox}
\makeatletter
\providecommand{\subtitle}[1]{% add subtitle to \maketitle
  \apptocmd{\@title}{\par {\large #1 \par}}{}{}
}
\makeatother

\author{A. Einstein, C.A de Coulomb \\ 
\emph{Institute for Nanofabrication, University College, Cambridge, MA
02139} \\ \emph{einstein@university.edu} \\
\vspace{3ex} \\ 
R. Fineman \\ 
\emph{HAL Laboratory of Quantum Studies, 137 Coupling Way, Pasadena, CA
91109}}

\date{\vspace{-3ex}}

\begin{document}
\maketitle


Abstracts should be supplied in Microsoft Word (.doc) or as an Adobe PDF
manuscript. Please use Times Roman, Helvetica or Arial as your font. 12
point or larger fonts should be used throughout, with the title
presented in 16 point or larger. Present the title and author list
centered, while the paragraphs of the abstract should be left justified.
Please ensure that the presenting author's name is \uline{underlined} in
the author list. Please also provide a contact e-mail.

Use margins of at least 1.5 inch (3.81 cm) on both the left and right
sides. Abstracts are limited to one page of text. References should
appear on the bottom of the first page.\footnote{References should
  follow standard APS style guidelines. For more, see
  http://publish.aps.org/STYLE/ms.html} A second, optional page with up
to four figures may also be included. Please be certain to make a
specific reference to the figures, such as to Figure 1. For the printed
book of Abstracts, images such as SEM micrographs will be printed with a
halftone process, and so you should bear in mind that images will be
reduced in quality when printed. Do not submit halftone images, since a
double halftoning processes will result in poorer quality.

You MUST use our on-line abstract submission system at
\href{http://www.eipbn.org}{www.eipbn.org}. Be sure to use the
``invited'' option for submission if your paper has been invited. Faxed
and e-mailed abstracts are not acceptable, but please contact the
Program Chair if you find the abstract submission is not working. Late
submissions may be considered, but post-deadline papers will be accepted
only if they report truly outstanding results. Abstracts will be peer
reviewed according to originality of the work, specific results
described, and potential impact. Please prepare abstracts carefully and
describe accomplishments specifically.

Please bear in mind that your abstract will be formatted for US standard
8.5x11 inch ``US Letter'' paper. Formatting your abstract for the longer
European A4 paper may result in lost text or cropped images. Please
ensure that any text in your figures is large enough to read.

Thank you very much for participating in the EIPBN conference.

\begin{figure}
\centering
\includegraphics[width=.6\textwidth]{figure.png}
\caption{\emph{The Debutantes' Ball:} At the debutantes' ball, young
nano-women gather in their finery at the edge of the stage to weep
because none of the nano-boys will dance with them. The micrograph shows
an array of 50 nm wide posts with a periodicity of 100 nm. The posts
consist of PMMA on top of an antireflection coating. The substrate
consists of a 250 nm thick layer of silicon nitride on silicon. Winner
of the 1995 ``Most Bizarre Micrograph'' prize. Submitted by Tim Savas,
Massachusetts Institute of Technology.}
\end{figure}

\end{document}


\end{document}